% Capitulo: Nuvem, big data e o modern data stack
\chapter{Nuvem, big data e o modern data stack}
\label{cap:nuvem}

Enquanto o self-service levava a análise para a mesa do analista, os bastidores
mudavam de escala. A internet, os celulares e os sensores passaram a gerar dados num
volume que os armazéns tradicionais não davam conta de guardar, e muito menos de
processar dentro de uma janela noturna. Este capítulo percorre as três ondas que vieram
em seguida: o big data, os armazéns de dados na nuvem e a troca do ETL pelo ELT, que deu
origem ao chamado modern data stack. E termina com uma constatação que vai acompanhar o
resto do livro: muitas das ideias novas são ideias antigas com outra roupa.

\section{Big data e data lakes}
\label{sec:big-data}

Em 2001, Doug Laney, analista do META Group, descreveu os desafios de gestão de dados em
três dimensões: volume, velocidade e variedade \parencite{laney2001}. Os \enquote{três
Vs} viraram, anos depois, a definição informal de big data: dados grandes demais, rápidos
demais ou heterogêneos demais para as ferramentas da época.

A resposta técnica veio das empresas de internet, que enfrentavam o problema antes de
todo mundo. O Google publicou em 2003 a arquitetura do seu sistema de arquivos
distribuído, que espalhava dados por milhares de máquinas baratas
\parencite{ghemawat2003}, e em 2004 o MapReduce, um modelo para processar esses dados
em paralelo dividindo o trabalho em duas etapas simples \parencite{dean2004}. Um projeto
de código aberto inspirado nesses dois artigos, o Hadoop, ganhou força a partir de 2006 e
se tornou sinônimo de big data na década seguinte.

A promessa do Hadoop era guardar tudo, em qualquer formato, a um custo baixo, e decidir
depois o que fazer com os dados. Em 2010, James Dixon cunhou o termo \emph{data lake}
para essa ideia: em vez de um armazém em que a água chega engarrafada e pronta para
consumo, um lago onde ela chega no estado natural e cada um retira o que precisa
\parencite{dixon2010}.

A metáfora era boa, e a prática nem sempre acompanhou. Muitos lagos viraram pântanos:
arquivos sem documentação, sem dono, sem garantia de qualidade, que ninguém tinha
coragem de apagar e ninguém conseguia usar. E a tecnologia exigia especialistas raros,
escrevendo código de baixo nível para tarefas que, num banco relacional, seriam uma
consulta SQL de poucas linhas. Não por acaso, uma das primeiras coisas que o ecossistema
Hadoop ganhou foram ferramentas para consultar os dados com SQL.

O big data deixou duas heranças duradouras. A primeira é a ideia de separar o
armazenamento barato e flexível do processamento, que vai reaparecer na nuvem e no
lakehouse. A segunda é a lição de que guardar tudo não substitui modelar e documentar. O
dado sem contexto, por mais barato que seja guardá-lo, continua sem valor.

\section{O armazém vai para a nuvem}
\label{sec:dw-nuvem}

A segunda onda trouxe o data warehouse de volta, agora como serviço. Em vez de comprar
servidores e licenças, a empresa passou a alugar capacidade de um provedor de nuvem e a
pagar pelo uso.

Três sistemas marcaram essa transição, e todos foram descritos em artigos técnicos que
valem a leitura. O Dremel, do Google, mostrou que era possível consultar bilhões de
linhas em segundos usando armazenamento em colunas e processamento distribuído
\parencite{melnik2010}; ele está por trás de um dos primeiros armazéns oferecidos como
serviço. A Amazon defendeu um armazém de dados mais simples de contratar e operar, que
qualquer time pudesse ligar em minutos \parencite{gupta2015}. E o Snowflake levou ao
extremo a separação entre armazenamento e processamento: os dados ficam num lugar, e
quantos \enquote{motores} de processamento forem necessários podem ser ligados e
desligados sobre eles, de forma independente \parencite{dageville2016}.

Essa separação mudou a economia do armazém de dados. No modelo antigo, a capacidade era
comprada para o pico e ficava ociosa a maior parte do tempo. No modelo novo, o custo
acompanha o uso: uma consulta pesada custa mais, um processo que roda a cada hora custa
mais do que um que roda uma vez por dia. É uma vantagem enorme para quem começa pequeno,
e uma armadilha para quem não acompanha o consumo. Essa mudança, da licença fixa para o
pagamento por uso, foi o primeiro passo de uma tendência que a IA acelerou e que o
capítulo~\ref{cap:custo-consumo} trata em detalhe.

\section{De ETL para ELT}
\label{sec:elt}

Com processamento barato e elástico dentro do próprio armazém, uma inversão passou a
fazer sentido. Em vez de transformar os dados num servidor intermediário antes de
carregá-los (ETL), a empresa carrega os dados brutos no armazém e transforma lá dentro,
com SQL (ELT). A janela noturna deixa de ser um gargalo, o dado bruto fica disponível
para reprocessar quando uma regra muda, e a transformação passa a ser escrita na
linguagem que os analistas já conhecem.

Em meados dos anos 2010 surgiram ferramentas para organizar essa transformação como
código: modelos em SQL, versionados, testados e documentados, com dependências entre si
resolvidas automaticamente. A mais conhecida delas, o dbt, ajudou a consolidar uma forma
de trabalhar que trouxe práticas de engenharia de software para dentro do time de dados.
A Parte~\ref{parte:devops} do livro é, em boa medida, sobre essas práticas.

Com a nova forma de trabalhar veio um papel novo. Em 2019, Michael Kaminsky descreveu o
\emph{analytics engineer}: alguém entre o engenheiro de dados e o analista, que domina SQL
e modelagem, aplica práticas de engenharia de software e entrega conjuntos de dados
limpos e bem documentados para os analistas usarem \parencite{kaminsky2019}. A
comunidade em torno do dbt adotou e popularizou o termo. O
capítulo~\ref{cap:mapas} compara esse papel com o do analista e o do engenheiro de
dados.

Os projetos de analytics engineering costumam se organizar em camadas: uma camada de
preparação, que limpa e padroniza o dado bruto de cada fonte; uma camada intermediária,
que junta e aplica regras de negócio quando necessário; e uma camada final, com os
modelos prontos para consumo, muitas vezes dimensionais. Quem leu o capítulo anterior
reconhece aqui ecos de Inmon e de Kimball. O capítulo~\ref{cap:medalhao} mostra como
essa organização em camadas ganhou um nome próprio, a arquitetura medalhão.

Ao conjunto de ferramentas dessa época (ingestão gerenciada, armazém na nuvem,
transformação como código, BI por cima) deu-se o nome de \emph{modern data stack}. O
nome é de marketing, mas a mudança foi real: um time pequeno passou a conseguir montar,
em semanas, uma plataforma que antes exigia um projeto de anos.

\rascunho{
  Relato: a experiência de montar ou migrar uma plataforma nesse modelo (armazém na
  nuvem, transformação em camadas, BI por cima), com o que deu certo e o que surpreendeu,
  inclusive no custo. Sem nomes de empresa e com números só em ordem de grandeza.
}

\section{Ideias antigas, nomes novos}
\label{sec:nomes-novos}

A área de dados tem o hábito de redescobrir as próprias ideias. O lakehouse, proposto em
2021, junta o armazenamento aberto e barato do data lake com as garantias de um data
warehouse: transações, esquemas, desempenho de consulta \parencite{armbrust2021}. É uma
evolução técnica real, e também uma volta ao ponto em que o data warehouse já estava,
agora sobre outra base.

A camada semântica, em que cada métrica é definida uma vez e reaproveitada por todas as
ferramentas, voltou à moda nos últimos anos. As suítes corporativas de BI dos anos 1990
já tinham camadas desse tipo. A organização em camadas de refinamento lembra a separação
entre a área de preparação e a área de apresentação de Kimball. E o
\enquote{dado como produto} lembra muito o data mart com dono.

Reconhecer isso não é desmerecer o novo. Cada volta acontece sobre uma tecnologia
melhor, mais barata e mais acessível, e por isso chega a mais gente. Mas reconhecer o
padrão ajuda a separar o que é de fato novo do que é rótulo, e a aproveitar décadas de
experiência acumulada em vez de reaprender tudo do zero. O próximo capítulo trata da
onda mais recente, e talvez da primeira em muito tempo que muda não só a infraestrutura,
mas também o trabalho de quem analisa os dados: a chegada da IA generativa.

\begin{leituras}
  \leitura{dageville2016}{O artigo que descreve a separação entre armazenamento e
    processamento no Snowflake. Ajuda a entender por que o custo da nuvem se comporta como
    se comporta.}
  \leitura{dean2004}{O artigo do MapReduce, base do Hadoop. Curto, claro e importante
    para entender a década do big data.}
  \leitura{kaminsky2019}{O texto que descreveu o papel do analytics engineer, ainda uma
    boa leitura para quem quer entender onde ele se encaixa entre análise e engenharia.}
  \leitura{reis2022}{Uma visão atual e independente de ferramenta do ciclo de vida da
    engenharia de dados, da geração ao consumo.}
\end{leituras}
